% Это основная команда, с которой начинается любой \LaTeX-файл. Она отвечает за тип документа, с которым связаны основные правил оформления текста.
\documentclass{article}[a4paper, 14pt]

% Здесь идет преамбула документа, тут пишутся команды, которые настраивают LaTeX окружение, подключаете внешние пакеты, определяете свои команды и окружения. В данном случае я это делаю в отдельных файлах, а тут подключаю эти файлы.

% Здесь я подключаю разные стилевые пакеты. Например возможности набирать особые символы или возможность компилировать русский текст. Подробное описание внутри.
\usepackage{packages}

% Здесь я определяю разные окружения, например, теоремы, определения, замечания и так далее. У этих окружений разные стили оформления, кроме того, эти окружения могут быть нумерованными или нет. Все подробно объяснено внутри.
\usepackage{environments}

% Здесь я определяю разные команды, которых нет в LaTeX, но мне нужны, например, команда \tr для обозначения следа матрицы. Или я переопределяю LaTeX команды, которые работают не так, как мне хотелось бы. Типичный пример мнимая и вещественная часть комплексного числа \Im, \Re. В оригинале они выглядят не так, как мы привыкли. Кроме того, \Im еще используется и для обозначения образа линейного отображения. Подробнее описано внутри.
\usepackage{commands}

% Пакет для титульника проекта
\usepackage{titlepage}

% Пакеты для отображения картинок
\usepackage{graphicx}
\usepackage{float}
\usepackage{wrapfig}

%
%\usepackage[unicode=true, colorlinks=true, linkcolor=blue, urlcolor=blue]{hyperref}


% Здесь задаем параметры титульной страницы
\setUDK{192.168.1.1}
% Выбрать одно из двух
%\setToResearch
%\setToProgram
\setToPractice

\setTitle{Мобильное приложение для отслеживания привычек}


\setStageOne

%сюда можно воткнуть картинку подписи
\setStudentSgn{}

\providecommand{\studentCompleted}{Выполнил студент}
\providecommand{\student}{TestStudent}
\providecommand{\group}{404}

\setStudent{\student}
\setGroup{\group}

\setStudentDate{20.10.2022}
\setAdvisor{Колесниченко Елена Юрьевна}
\setAdvisorTitle{департамент больших данных и информационного поиска ФКН ВШЭ, доцент}
\setAdvisorAffiliation{ООО <<Яндекс. Технологии>>}
\setAdvisorDate{24.10.2022}
\setGrade{--}
%сюда можно воткнуть картинку подписи
\setAdvisorSgn{\scalebox{0.4}{\includegraphics{Utils/advisor_sign}}}
\setYear{2022}

\setStudentCompleted{\studentCompleted}

\linespread{1.2}
\urlstyle{same}

\bibliographystyle{Bibliography/gost-numeric.bbx}
\usepackage[parentracker=true,
backend=biber,
hyperref=true,
bibencoding=utf8,
language=auto,
autolang=other,
citestyle=gost-numeric,
defernumbers=true,
bibstyle=gost-numeric,
sorting=ntvy,
maxbibnames=99,
]{biblatex}

\addbibresource{Bibliography/sources.bib}

\setlength{\parskip}{6pt}
\setlength{\parindent}{10pt}

% С этого момента начинается текст документа
\begin{document}

% Эта команда создает титульную страницу
\makeTitlePage

% Здесь будет автоматически генерироваться содержание документа
\tableofcontents


\newpage
\input{Parts/intro}

\newpage
%!TEX root = ../main

\section{Large Scale Distributed Deep Networks}
%\indent\setlength{\parindent}{1em} 
Начнем изучение темы распределённого обучения моделей со статьи “Large Scale Distributed Deep Networks\cite{beginning}, в которой делается большой шаг в сторону распределенного обучения нейронных сетей. 

%\indent\setlength{\parindent}{1em} 
Нейронная сеть состоит из нескольких слоев нейронов и синапсов, связей между нейронами. Каждый нейрон принимает на вход какие-то небольшие данные, производит над ними работу и передает полученный результат по синапсу следующему нейрону. У каждого синапса есть вес, который влияет на то, какой вклад внесет результат вычислений одного нейрона на результат вычислений следующего. 


\begin{figure}[h]%current location
	\centering
	\scalebox{0.5}{\includegraphics{Parts/images/neural_network.jpg}}
	\caption{Визуализация нейронной сети}
	\label{framework} %framework,fig1
\end{figure}

%\indent\setlength{\parindent}{1em} 
Веса синапсов - это одни из параметров нейронной сети и в задачу обучения сети входит подбор весов синапсов так, чтобы минимизировать ошибку модели. Таким образом, подбор весов нейронной сети сводится к задаче минимизации функции ошибки.

%\indent\setlength{\parindent}{1em} 
Основной вклад, который внесла статья в развитие распределённого обучения моделей, заключается в том, что авторы разработали два распределенных метода оптимизации и провели серии экспериментов, показывающих эффективность этих моделей по сравнению с нераспределёнными, работающими на CPU или на соптимизированном GPU.

%\indent\setlength{\parindent}{1em} 
В статье авторы описывают два распределенных метода оптимизации, которые можно применять для поиска весов нейронной сети: Downpour SGD и Sandblaster L-BFGS. Оба метода были разработаны для фреймворка DistBelief, поддерживающего параллелизм моделей, когда модель делится на части и каждая часть обучается на одном и том же датасете, и параллелизм данных, когда копия модели обучается на разных поднаборах датасета. 

\subsection{Downpour SGD}
%\indent\setlength{\parindent}{1em}

\begin{wrapfigure}[8]{r}{0.4\textwidth}
	\vspace{-2cm}
	\begin{center}
		\includegraphics[width=0.9\linewidth]{Parts/images/Downpour_SGD.jpg}
	\end{center}
	\caption{Визуализация Downpour SGD}
\end{wrapfigure}

Downpour SGD - асинхронная вариация градиентного спуска SGD. В основе распределенности этой модели лежит параллелизм данных. Несколько копий модели запускаются на разных репликах, каждой реплике отведен поднабор параметров, который она будет менять во время обучения. Модели обмениваются обновленными параметрами с помощью центрального сервера параметров, который хранит текущее состояние всех параметров, которые используются на разных репликах.

% \begin{figure}[h]%current location
% 	\centering
% 	\scalebox{0.8}{\includegraphics{Parts/images/Downpour_SGD.jpg}}
% 	\caption{Визуализация Downpour SGD}
% 	\label{framework} %framework,fig1
% \end{figure}

%\indent\setlength{\parindent}{1em} 
Преимуществом Downpour SGD по сравнению с синхронной распределенной вариацией SGD является его устойчивость к отказам машин, поскольку каждая реплика индивидуально выполняет свою работу и по завершении не ждет остальных, чтобы обновить значения параметров. Основной проблемой асинхронного подхода является невозможность поддержания консистентности значений параметров для каждой реплики. Поскольку каждая реплика действует независимо, то нет никакой гарантии, что в любой момент времени к параметрам на разных машинах было применено одинаковое количество обновлений или что обновления были применены в одном и том же порядке.

\subsection{Sandblaster L-BFGS}
\begin{wrapfigure}[14]{r}{0.5\textwidth}
	\begin{center}
		\includegraphics{Parts/images/Sandblaster_L-BFGS.jpg}
	\end{center}
	\caption{Визуализация Sandblaster L-BFGS}
\end{wrapfigure}

%\indent\setlength{\parindent}{1em} 
Авторы рассматривают распределенное обучение модели L-BFGS, заключающееся в параллелизме данных.  

%\indent\setlength{\parindent}{1em} 
Ключевые идеи распределенности для алгоритма L-BFGS, которые предлагают авторы статьи:
\begin{itemize}
  \item распределенное хранение параметров 
  \item выделение процесса координатора, который не имеет доступа к параметрам модели и выступает в роли оркестратора всех реплик, выполняющих обучение
\end{itemize}

%\indent\setlength{\parindent}{1em} 
Координатор назначает каждой из реплик модели небольшую часть работы, которую она может выполнить независимо, и каждый раз, когда какая-то реплика заканчивает свою работу, координатор назначает ей новую часть. Такой подход позволяет хорошо использовать быстрые машины и исключает ожидания самой медленной реплики, чтобы продолжить процесс. Результаты работы, а также кэшируемая информация сохраняются локально на каждой реплике и не отсылаются в центральный сервер, что позволяет запускать обучение достаточно больших моделей без оверхеда на отправку и сбор всех параметров. 


\subsection{Достигнутые результаты и проверка эффективности}
%\indent\setlength{\parindent}{1em} 

% \begin{wrapfigure}[14]{r}{0.7\textwidth}
% 	\begin{center}
% 		\includegraphics[width=0.9\linewidth]{Parts/images/time_accuracy.jpg}
% 	\end{center}
% 	\caption{Зависимость точности модели от времени обучения}
% \end{wrapfigure}

Авторы статьи применяют описанные методы оптимизации для двух задач глубинного обучения: распознавание объектов на изображениях и акустическая обработка для распознавания речи.
Ниже приведены графики для задачи обработки речи. Авторы сравнивают три оптимизационных модели, которые тренируются распределённо: Downpour SGD (фиолетовая линия), Downpour SGD w Adagrad (красная линия), Sandblaster L-BFGS (зеленая линия) с двумя моделями, одна из которых была обучена на центральном процессоре одной реплики (черная линия), а другая - на оптимизированном графическом процессоре (розовая пунктирная линия). Опираясь на графики можно сказать, что модели, обученные распределенным способом, намного быстрее увеличивают точность решения задачи. 

\begin{figure}[H]%current location
	\centering
	\scalebox{0.9}{\includegraphics{Parts/images/time_accuracy.jpg}}
	\caption{Зависимость точности модели от времени обучения}
	\label{framework} %framework,fig1
\end{figure}

%\indent\setlength{\parindent}{1em} 
Весьма значительным результатом является то, что в статье показали возможность подхода распределенного обучения масштабироваться до моделей, которые намного больше, чем можно удобно разместить на одной машине, тем более на одном графическом процессоре. Авторы применили разработанный метод обучения модели оптимизации Downpour SGD к нейронной сети с 1.7 миллиардами параметров, которая решает задачу классификации изображения, и получили точность классификации 15\% на кросс валидации, что является относительным улучшением более чем на 60\% по сравнению с лучшей производительностью в этой задаче классификации. 


\newpage
\input{Parts/all_reduce}

\newpage
%!TEX root = ../main

\newpage
\section{Model parallel \& data parallel}
%\indent\setlength{\parindent}{1em} 
Нейронные сети, в отличие от классических моделей машинного обучения, требуют большого времени на обучение, иногда обучение может достигать нескольких недель или месяцев. Естественной идеей решения такой проблемы является попытка создания асинхронного обучения. С такой целью были созданы 2 основных подхода: data-parallelism и model-parallelism. Однако оба эти способа обладают недостатками, не позволяющими реализовать весь потенциал распределенных вычислений. В статье "PipeDream: Fast and Efficient Pipeline Parallel DNN Training"\cite{pipedream} описаны основные проблемы данных подходов и предложена система PipeDream, комбинирующая эти методы и тем самым оптимизирующая их.

\subsection{Data parallelism}
Идея Data parallelism заключается в разбиении исходных данных на части, каждая из которых будет использована на отдельной GPU. На каждой GPU находится своя полная копия Нейронной Сети, и веса, получаемые в ходе обучения, синхронизируются между GPU раз в несколько эпох.


\begin{figure}[h]%current location
	\centering
	\scalebox{0.8}{\includegraphics{Parts/images/data-parallel.jpg}}
	\caption{Пример data-parallel на 4 GPU}
	\label{data_parallel} %framework,fig1
\end{figure}


%\indent\setlength{\parindent}{1em} 
Однако главная проблема data-parallel заключается в долгой синхронизации данных между процессами. Были проведены тесты на известных нейронных сетях, используя 3 разные поколения NVIDIA GPU -- Kepler (K80), Pascal (Titan X) и Volta (V100):

\begin{figure}[h]%current location
	\centering
	\scalebox{0.7}{\includegraphics{Parts/images/communication-overhead.jpg}}
	\caption{Время на синхронизацию в \% от всего времени обучения DNN}
	\label{communication_overhead} %framework,fig1
\end{figure}

%\indent\setlength{\parindent}{1em} 
Как видно на данном графике \ref{communication_overhead}, время, затрачиваемое на синхронизацию , может достигать до 90\% от всего времени обучения DNN.


\subsection{Model parallelism}
Идея Model Parallelism заключается в разбиении самой модели на несколько последовательных частей, каждая из которых будет вычисляться на отдельной GPU. Данный подход часто используется в Машинном Обучении (ML). Однако для DNN наивный подход model-parallel имеет огромный недостаток: каждая следующая часть слоёв должна дожидаться завершения работы предыдущей.

\begin{figure}[h]%current location
	\centering
	\scalebox{1}{\includegraphics{Parts/images/model-parallel.jpg}}
	\caption{Процесс наивного model-parallel на 4 GPU. \\Числа обозначают номер батча.}
	\label{framework} %framework,fig1
\end{figure}

%\indent\setlength{\parindent}{1em} 
Из-за такого подхода наши GPU значительную часть времени находятся в режиме ожидания, и как такового прироста в производительности нет.

\subsection{PipeDream}
%\indent\setlength{\parindent}{1em}  
PipeDream использует идею model-parallel, разделя DNN на последовательные части по несколько слоев Нейронной Сети, и при этом реализуют возможность параллельно выполнения прямого и обратного проходов для разных батчей, тем самым минимизируя простои GPU.

\begin{figure}[h]%current location
	\centering
	\scalebox{1}{\includegraphics{Parts/images/pipelined_model_example.jpg}}
	\caption{Процесс обучения PipeDream, объединяющий data-parallel и model-parallel}
	\label{framework} %framework,fig1
\end{figure}

\newpage
\subsubsection{Распределение слоев между GPU}

\begin{wrapfigure}[14]{r}{0.5\textwidth}
    \vspace{-2cm}
	\begin{center}
		\includegraphics[width=\linewidth]{Parts/images/pipedream_process.jpg}
	\end{center}
	\caption{Автоматизированный процесс деления DNN на части}
    \label{fig:transformer1}
\end{wrapfigure}

%\indent\setlength{\parindent}{1em}  
Для минимизации времени работы GPU в режиме ожидания важно правильно поделить DNN на части. Главная задача при разделении заключается в том, что каждая отдельная компонента PipeDream выполнялась одинаковое количество времени.

Для этого сначала производится замер времени работы DNN, производя вычисления 1000 батчей на одной из GPU. После этого при помощи алгоритма динамического программирования находится оптимальное разбиение слоёв модели на части, каждая из который будет вычисляться одинаковое количество времени (тогда мы можем считать, что каждая часть выполняется ровно одну единицу времени).

% \begin{figure}[h]%current location
% 	\centering
% 	\scalebox{0.7}{\includegraphics{Parts/images/pipedream_process.jpg}}
% 	\caption{Автоматизированный процесс деления DNN на части}
% 	\label{framework} %framework,fig1
% \end{figure}

\subsubsection{Последовательность вычисления}
%\indent\setlength{\parindent}{1em}
В отличие от обычных реализаций pipeline, в которых одновременно вычисления происходят только либо в прямом, либо в обратном порядке, в PipeDream вычисления происходят в обе стороны одновременно. То есть у каждой GPU в каждую единицу времени есть выбор: либо делать прямой проход графа вычислений, либо обратный. 

\begin{figure}[h]%current location
	\centering
	\scalebox{1}{\includegraphics{Parts/images/pipedream_work_schedule.jpg}}
	\caption{Автоматизированный процесс деления DNN на части}
	\label{framework} %framework,fig1
\end{figure}

%\indent\setlength{\parindent}{1em}  
В начальной стадии использнения PipeDream, первая часть получает на вход оптимальное количество батчей, последовательно после вычисления передавая следующей части (а значит и следующей GPU) результат вычисления. 

%\indent\setlength{\parindent}{1em}  
В основной стадии работы PipeDream на одной GPU чередуются прямой и обратный проходы. Такой механизм авторы статьи назвали one-forward-one-backward (1F1B).

\subsubsection{Эффективность обучения}
%\indent\setlength{\parindent}{1em} 
Как мы можем видеть из примера из последней иллюстрации: прямой проход для батча №5 для первой GPU проходит после обновления весов после обратного прохода батча №1. При этом обратный проход батча №5 проходит после оборатных проходов батчей №2, №3 и №4. Это может негативно влиять на сходимость.

%\indent\setlength{\parindent}{1em}
\textbf{Weight Stashing:} в PipeDream реализована идея сохранения весов для всех батчей, которые в данный момент времени вычисляются в модели. Таким образом, когда батчу нужно пройти прямой проход, то используются последняя версия весов. После этого сохраняется версия весов, которая вспоследствии будет использована при обратном проходе того же батча. Такой метод авторы статьи называют Weight Stashing. Благодаря ему для каждой части модели для каждого батча используются подходящие веса, как для прямого, так и для обратного проходов. Тем самым не будет нарушена сходимость графа вычислений.

\subsection{Заключение}
%\indent\setlength{\parindent}{1em} 
PipeDream существенно превосходит метод data-parallel, поскольку каждая GPU обменивается только частью параметров, тем самым ускоряя их синхронизацию:

\begin{figure}[h]%current location
	\centering
	\scalebox{0.7}{\includegraphics{Parts/images/pipedream_tests.jpg}}
	\caption{Сравнение PipeDream с data-parallelism (BSP)}
	\label{framework} %framework,fig1
\end{figure}

%\indent\setlength{\parindent}{1em} 
По результатам экспериментов для моделей VGG16 и S2VT время, затраченное на синхронизацию, уменьшилось на 90\%, при этом в целом модель стала в среднем обучаться в 3 раза быстрее.

\newpage
\input{Parts/tensor_parallelism}

\newpage
%!TEX root = ../main

\section{Заключение}

Исходя из содержания, становится ясно, что распределенные системы обучения нейронных сетей ползволяют получить ранее недостигаемые результаты. Так как нейронные сети могут сильно различаться по количеству параметров, по типу слоёв и тому подобному, каждая представленная система лучше справляется именно с теми задачами, для которых она разработана. При этом, из общего, всем таким системам свойственно добавление накладных расходов и слабая консистентность данных, но в конечном счете увеличение самого главного — скорости достижения моделью целевой точности, что может быть не просто удобным механизмом для быстрого обучения, но и необходимостью, в случаях когда модель не может обучиться на одной машине за разумное время. Таким образом, использование распределённых систем становится неотъемлемым компонентом в машинном обучении, позволяющим справляться с ранее нерешаемыми задачами.


\newpage 
%!TEX root = ../main

\section{Распределение обязанностей в команде}

\begin{itemize}
    \item Иван Ершов -- анализ статьи \cite{pipedream}
    \item Дарья Жукова -- анализ статьи \cite{beginning}
    \item Ксения Петренко -- анализ статьи \cite{shoeybi2019megatron}
    \item Тагир Хамитов -- анализ статьи \cite{DBLP}
    \item Артём Макогон -- сборка, оформление \href{https://github.com/makogon2907/HSE-Practice-Fall-2022}{отчета}
    \item Дмитрий Брусенин -- сборка, оформление \href{https://github.com/makogon2907/HSE-Practice-Fall-2022}{отчета}
\end{itemize}

\newpage
\addcontentsline{toc}{section}{Список литературы}
\printbibliography

% \addcontentsline{toc}{section}{Список иллюстраций}
% \listoffigures




\end{document}
% Начиная с этого момента весь текст LaTeX игнорирует, можете вставлять любую абракадабру.
